% Introduction to Part VII. Unnumbered text: no \label except the \phantomsection target.
\phantomsection\label{intro:VII}
Part~VII constructs one step of the small-field renormalization group and writes it exactly as a Gaussian integral around the background field. The step $T$ integrates over the configurations $V$ on $\T=\T^{(k)}$ whose one-step average is a given configuration $V'$ on $\T'=\T^{(k+1)}$ (Section~\ref{sec:rg-step}). In four dimensions a change of units does not change the coupling; the coupling of the next step comes from the fluctuation integral, and Part~X reads it off from the new generation of activities (Section~\ref{sec:rg-step}, paragraph ``The couplings''). Part~VII proves the exact representation and the linear estimates on which Parts~VIII--X rest (Remark~\ref{x:34.11}). The gauge group is semisimple from here on.

\paragraph{What this Part proves.} The main theorem is Theorem~\ref{thm:K}, assembled from Theorems~\ref{thm:step} and~\ref{thm:step-coercive}, Corollary~\ref{cor:covariance}, Lemma~\ref{lem:gf-step-invariance} and Proposition~\ref{prop:symmetry}. The anatomy of the step is as follows.
\begin{itemize}
\item \emph{Gauge fixing.} Ba\l aban's exponential gauge-fixing function $\mathrm{gf}_W$ fixes the invariance under the gauge transformations of $\T$ that are trivial on $\T'$, in the averaged contour variables relative to the base $W=U'^{(k)}$, the $k$-step average of the background field $U'$ of $V'$ (Definition~\ref{def:gauge-fixing}, Lemma~\ref{lem:composition}). By a Faddeev--Popov identity without determinant (Lemma~\ref{lem:exp-gauge}), the gauge-fixed step agrees with $T$ on invariant integrands supported on regular configurations; for $\rho\ge0$ the part removed by the cutoffs is nonnegative (Lemma~\ref{lem:gf-step-invariance}).
\item \emph{Elimination.} The constraint is solved on the middle bond $b_0(c)$ of each coarse bond $c$, with local Jacobians, and the fluctuation field lives on the set $\mathfrak F$ of the other bonds (Lemma~\ref{lem:eliminate}, Proposition~\ref{prop:exact-step}).
\item \emph{The exponent.} The coarse effective action $V\mapsto S(U_k(V))$ is analytic, its gradient is the Lagrange multiplier, and its derivatives decay along trees (Proposition~\ref{prop:coarse-action}). After composition with the eliminated field $\hat Z(Z)$ and addition of the gauge-fixing exponent $\Phi_W(Z)$, the quadratic part is $\frac12\Delta_k(Z,Z)$; the remainder is the sum of the third-order coarse-action interaction $\mathcal V(Z)$ and the block-local gauge-fixing remainder $G_3(Z)$ (Lemma~\ref{lem:step-quadratic}).
\item \emph{Coercivity.} The fluctuation operator satisfies $\frac14c_L\|Z\|^2\le\Delta_k(Z,Z)\le C_L\|Z\|^2$, and its inverse $C_k$ decays exponentially (Theorem~\ref{thm:step-coercive}, Corollary~\ref{cor:covariance}). The input is a Poincar\'e inequality for one step, in which the gauge-fixing term cannot be dropped (Lemma~\ref{lem:poincare-step}).
\item \emph{The representation.} For $k\ge1$, the gauge-fixed step $T^{\rm gf}(\chi^{\rm w}\rho)(V')$ equals an explicit constant times $e^{-g^{-2}S(U')}(\det\Delta_k)^{-1/2}\prod_c\mathfrak j_c$ times a Gaussian integral with covariance $C_k$ (Theorem~\ref{thm:step}), and $\log\det\Delta_k$ is a sum of quasi-local, gauge-invariant densities (Proposition~\ref{prop:det-densities}).
\end{itemize}
The curvature of the background field of a coarse configuration of small curvature is at most $(C_*+\delta)$ times that curvature, with $C_*$ independent of $L$ (Proposition~\ref{prop:sharp-curvature}, with Lemma~\ref{lem:uniformity}). So the small-field thresholds of consecutive steps are compatible for $L^2\ge8(C_*+1)$, provided $g_{k+1}\le2g_k$ and the fluctuation cutoff is small, as in condition~(Q2) of Remark~\ref{rem:parameters}. The step commutes with the lattice symmetries (Proposition~\ref{prop:symmetry}). The constants do not depend on $k$, $n$, $\eta$ or the coupling.

\paragraph{Imports and exports.} Part~VII uses the background fields and the variational problem (Theorems~\ref{thm:variational} and~\ref{thm:balaban}, Corollary~\ref{cor:analytic}), the averaging (Propositions~\ref{prop:Q-expansion} and~\ref{prop:linearized-averages}, Lemma~\ref{lem:avg-transport}), and the operators of Parts~II and~III at regular backgrounds. No proof uses a later result; the white-noise cutoff of the small-field step is Definition~\ref{def:small-field-step}, and the step from scale $0$ is Lemma~\ref{lem:first-step-formula}. Theorem~\ref{thm:step}, Corollary~\ref{cor:covariance}, Proposition~\ref{prop:det-densities}, Lemma~\ref{lem:eliminate} and Proposition~\ref{prop:exact-step} are the inputs of Parts~VIII--XI. Proposition~\ref{prop:hessian} goes to Parts~VIII, X and~XI, Proposition~\ref{prop:sharp-curvature} to Parts~VIII--X, Lemma~\ref{lem:gf-step-invariance} to Parts~X and~XI, and Remark~\ref{rem:one-loop-intrinsic} to Part~XI. All of Part~VII is on the main line.

\paragraph{Details for this Part.} The companion, chapter on Part~VII, expands each step of this Part, from the uniformity of the constants to Theorem~\ref{thm:step}, and the order of choice of its constants. A proof with an expansion there carries the marginal link ``Details''.
